\section{Neyman's confidence constructions: complete proofs}
\label{sec:annex-p15-neyman-confidence-constructions}

\begin{annexitem}{annex:p15.neyman.uniform.coverage}{thm:p15-neyman-uniform-coverage}{The three uniform systems and every section}
\begin{proof}
The joint density is \(\theta^{-2}\) on \((0,\theta)^2\)
and zero elsewhere. For \(A_1^{\rm U}(\theta)\), each coordinate
is in its specified interval with probability \(1-\beta\),
so independence gives \((1-\beta)^2=\alpha\).
Solving \(\beta t\leq x_i\leq t\) for both coordinates gives
exactly \(L\leq t\leq l/\beta\). This is empty when
\(l<\beta L\), a singleton at equality and the stated
closed interval otherwise.
The joint minimum/maximum density is \(2/\theta^2\) for
\(0<l<L<\theta\), since either coordinate may be the smaller.
Consequently the probability of an empty section is
\[
               \int_0^\theta\int_0^{\beta L}
                           \frac2{\theta^2}\,dl\,dL=\beta.
\]
Thus the nonempty-interval requirement genuinely fails on a
positive-probability set, while the general confidence-set
coverage remains exactly \(\alpha\).

For the sum, its convolution density at \(s\) is
\(\theta^{-2}\) times the length of
\((0,\theta)\cap(s-\theta,s)\). This gives the two
linear pieces in the statement and zero outside.
For a half-width \(\Delta\in[0,\theta]\), integrating
the two excluded triangles gives
\[
                    P_\theta(|S-\theta|\leq\Delta)
                         =1-(1-\Delta/\theta)^2.
\]
The choice \(\Delta=(1-q)\theta\), \(q=\sqrt{1-\alpha}\),
therefore gives coverage \(\alpha\).
The acceptance inequalities are \(qt\leq S\leq(2-q)t\);
their exact inverse is \([S/(2-q),S/q]\).
It is nonempty for every \(S>0\), because \(0<q<1\).

For \(n\) observations, independence gives
\(P_\theta(L_n\leq z)=\prod_iP_\theta(X_i\leq z)
=(z/\theta)^n\) on \(0\leq z\leq\theta\).
Thus
\[
                   P_\theta(q_n\theta\leq L_n\leq\theta)
                          =1-q_n^n=\alpha.
\]
Inversion yields \([L_n,L_n/q_n]\), always nonempty.
Every equality boundary used here is a probability-zero
set under every fixed continuous sampling law. The
declared closed inequalities nonetheless matter for
the exact set returned at each possible observation.
\end{proof}
\end{annexitem}

\begin{annexitem}{annex:p15.neyman.uniform.lengths}{thm:p15-neyman-uniform-lengths}{All lengths, source numbers and repetition probabilities}
\begin{proof}
Subtracting endpoints gives
\[
 D_2=S\left(\frac1q-\frac1{2-q}\right)
       =\frac{2(1-q)S}{q(2-q)},\qquad
 D_3=L(1/q-1).
\]
Since the observations are positive,
\(S/2\leq L<S\), with the first equality precisely
when they agree. Dividing the two positive lengths gives
\[
                         D_3/D_2=\frac{(2-q)L}{2S},
\]
and proves both bounds and their endpoint conditions.
Integration of a uniform density gives \(E_\theta S=\theta\).
Differentiating the maximum cumulative law gives density
\(2L/\theta^2\), so \(E_\theta L=2\theta/3\).
These expectations give the displayed expected lengths.
At \(\alpha=3/4\), \(q=1/2\), so the sum interval is
\([2S/3,2S]\), the maximum interval is \([L,2L]\),
and the numerical substitutions give every stated interval
and the ratio range \([3/8,3/4)\).

For trial \(j\), write its observations as
\(\theta_j(U_j,V_j)\), with independent uniform pairs
\((U_j,V_j)\) across trials. Each coverage condition
is invariant under this positive scale change.
Thus each procedure separately has independent Bernoulli
coverage indicators with parameter \(3/4\), regardless
of the deterministic positive sequence \(\theta_j\).
Their sum over 400 trials has the stated binomial law.
A frequency below \(0.69\) means at most 275 successes;
one above \(0.81\) means at least 325. This proves
the exact tail formula, which is approximately
\(0.0046303335\).

The indicators for the two procedures in one trial
are dependent. With scale one, the sum covers if
\(1/2\leq U+V\leq3/2\), and the maximum covers if
\(\max(U,V)\geq1/2\). Their joint success region has
area \(5/8\), obtained by removing from the sum's
area \(3/4\) the upper half of the square
\((0,1/2)^2\), of area \(1/8\).
Since \(5/8\ne(3/4)^2\), the indicators are not
independent. Across 400 independent pairs their
coverage counts have covariance
\(400[5/8-9/16]=25\), and are not independent either.
No conclusion about the realized entries of a particular
published random-number table has been used.
\end{proof}
\end{annexitem}

\begin{annexitem}{annex:p15.neyman.conditional.similarity}{thm:p15-neyman-conditional-similarity}{The exact conditional-factorization construction}
\begin{proof}
At a marginally relevant \(z\), division of the joint
density by \(m_{\theta,\eta}(z)>0\) gives the normalized
conditional density \(h_\theta(v\mid z)\), independent
of the nuisance value. More formally, for measurable
sets \(D,H\), integration of the joint density gives
\[
 P_{\theta,\eta}(T\in D,V\in H)
   =\int_Dm_{\theta,\eta}(z)
                       \left(\int_Hh_\theta(v\mid z)\,dv\right)dz.
\]
This is the conditional-probability identity. Values on
marginal null sets may be defined arbitrarily as probability
kernels without changing it. The stated coordinate
representation transfers it back to the original observations.

For an acceptance family with the required measurable
sections, product integration gives
\[
 P_{t,\eta}(V\in B_t(T))
  =\int m_{t,\eta}(z)
                 \left(\int_{B_t(z)}h_t(v\mid z)\,dv\right)dz
  =\alpha\int m_{t,\eta}(z)\,dz=\alpha.
\]
This proves the combination step for every nuisance value.
The calculation applies to any measurable sections;
it does not require them to be single intervals.
For a residual density on \(\mathbb R^r\), \(r\geq1\), let
\[
 H_\theta(z,y)=\int_{\{v_1\leq y\}}h_\theta(v\mid z)\,dv.
\]
Product integration makes this jointly measurable.
The conditional first-coordinate distribution has no atoms,
so its cumulative function is continuous and tends from zero
to one. Its value at the random first coordinate is uniform
on \((0,1)\): for \(0<c<1\), the upper endpoint of the
set of \(y\) with \(H_\theta(z,y)\leq c\) is finite,
has cumulative value \(c\) by continuity, and the probability
of that set is therefore \(c\). Flat stretches of the
cumulative function carry zero conditional probability;
the endpoint cases follow by limits.
For any \(b\in[0,1-\alpha]\), the section
\[
       B_{t,b}(z)=\{v:b\leq H_t(z,v_1)\leq b+\alpha\}
\]
is consequently measurable and has conditional probability
\(\alpha\). Distinct choices of \(b\) give distinct probability
coordinate intervals of positive symmetric-difference length,
so there are infinitely many choices.
It also identifies the required premise for selecting
their probabilities. For a conditional point mass,
nonrandomized sections can have probabilities only
zero or one, so arbitrary intermediate conditional
size would be impossible. In the triangular example
the conditional uniform law supplies sections of
every required size explicitly.
\end{proof}
\end{annexitem}

\begin{annexitem}{annex:p15.neyman.triangle.conditional}{thm:p15-neyman-triangle-conditional}{The full triangular distribution and nuisance calculation}
\begin{proof}
Write \(s=\theta_1\), \(\eta=\theta_2\), and \(\rho=s^3\eta\).
On its positive triangle the density is
\[
                    s^{-2}[2-\rho+3\rho(x_1/s)].
\]
For \(\rho\geq0\), its minimum over the closed unit
range of \(x_1/s\) is \(2-\rho\geq0\); at \(\rho=2\)
it is strictly positive whenever \(x_1>0\).
For \(\rho<0\), its minimum is \(2+2\rho>0\).
Thus it is positive in the open triangle and nonnegative
everywhere. Its integral is
\[
 \int_0^s\left(\frac2{s^2}-s\eta+3\eta x\right)(s-x)\,dx
                  =1-\frac{\eta s^3}{2}+\frac{\eta s^3}{2}=1.
\]
This also proves the stated marginal by integrating
over \(0<x_2<s-x_1\). Dividing by that positive
marginal on \(0<x_1<s\) gives \(1/(s-x_1)\),
the conditional uniform density, with no nuisance
parameter. At \(x_1=s\) the conditional expression
is not used; the marginal gives that event probability zero.

For the radial cutoff, integrate over the smaller
triangle \(x_1+x_2<a s\):
\[
 \int_0^{as}\left(\frac2{s^2}-s\eta+3\eta x\right)(as-x)\,dx
                    =a^2-\frac12a^2(1-a)\rho.
\]
Subtracting from one proves the acceptance probability.
For \(0<a<1\), the coefficient of \(\rho\) is
strictly positive, while the allowed nuisance range
contains more than one value. The probability
therefore cannot stay equal to one prescribed
coefficient for all nuisances. This is the exact
failure of the radial proposal.
\end{proof}
\end{annexitem}

\begin{annexitem}{annex:p15.neyman.triangle.coverage}{thm:p15-neyman-triangle-coverage}{The fixed conditional placements and their precise length class}
\begin{proof}
Given \(X_1=x_1\in(0,t)\), the null conditional
law is uniform on an interval of length \(t-x_1\).
The selected interval has length \(\alpha(t-x_1)\),
because its endpoint coefficients differ by \(\alpha\).
The conditions \(b>0\), \(b+\alpha\leq1\)
place it within the conditional support.
The conditional-similarity theorem therefore gives
probability \(\alpha\) for every valid nuisance value.

Solving the two inequalities for \(t\) gives
\[
                 x_1+\frac{x_2}{b+\alpha}
                      \leq t\leq x_1+\frac{x_2}{b}.
\]
These endpoints exceed \(x_1\) and are ordered
strictly since \(x_2>0\) and \(\alpha>0\).
The set is therefore nonempty for every sample.
Subtracting the endpoints yields
\(\alpha x_2/[b(b+\alpha)]\).
For \(b>0\), its denominator strictly increases,
so its length strictly decreases with \(b\).
The unique largest allowed \(b=1-\alpha\)
therefore gives the pointwise length minimum
within this whole fixed-placement family.
\end{proof}
\end{annexitem}

\begin{annexitem}{annex:p15.neyman.switched.intervals}{thm:p15-neyman-switched-intervals}{Every branch, switch boundary and crossing comparison}
\begin{proof}
For fixed \(t\), almost every conditional fiber
uses either its upper interval of relative length
\(\alpha\) or its lower interval of relative length
\(\alpha\). The exceptional fiber \(x_1=t/2\)
has marginal probability zero under every density.
Giving it the full conditional support, and closing
the endpoint inequalities on the other fibers,
therefore preserves total acceptance probability
\(\alpha\) for every nuisance value.

Fix positive data \(u,v\), and put \(r=v/u\).
The branch \(u<t/2\) gives
\[
              t>2u,\qquad u+v\leq t\leq u+v/(1-\alpha).
\]
The branch \(t/2<u<t\) gives
\[
              u+v/\alpha\leq t<2u.
\]
At \(t=2u\), the completed fiber accepts exactly
when \(v\leq u\), that is, \(r\leq1\).

If \(r>1\), only the upper branch is possible;
its closed interval is \([u+v,u+v/(1-\alpha)]\).
The formulas in the statement give these endpoints
because \(u+v>2u\).
If \(0<r\leq1\), the point \(2u\) is included.
The lower branch extends left to \(u+v/\alpha\)
exactly when \(r<\alpha\), and otherwise contributes
no point below \(2u\). The upper branch extends
right to \(u+v/(1-\alpha)\) exactly when
\(r>1-\alpha\), and otherwise contributes no
point above \(2u\). Equality at either
threshold gives only the already included
point \(2u\).
Thus all possibilities form one closed interval:
its lower endpoint is
\(\max\{u+v,u+\min(v/\alpha,u)\}\),
and its upper endpoint is the stated maximum.
When neither branch extends, this interval
is the singleton \(\{2u\}\).
This proves inversion and nonemptiness for every
\(\alpha\), including \(\alpha<1/2\).
The endpoints are continuous min/max expressions
in the positive data, so measurability is explicit.

At \(r=\alpha\), the lower endpoint is \(2u\),
and the upper is
\(\max\{2u,u/(1-\alpha)\}\).
Their difference is the displayed positive-part
formula. If \(\alpha>1/2\), subtraction from
the fixed-placement length gives
\[
 \frac{\alpha^2-(2\alpha-1)}{1-\alpha}u=(1-\alpha)u>0.
\]
If \(\alpha\leq1/2\), the switched length
is zero while the fixed-placement length
is positive. Thus it is strictly shorter
on that ray in every case.

The number
\(\alpha(1-\alpha)/(1-\alpha+\alpha^2)\)
is strictly smaller than both \(\alpha\)
and \(1-\alpha\): the two comparisons
reduce to positivity of \(\alpha^2\)
and of \((1-\alpha)^2\).
For \(r\) below it, only the lower
branch extends and the switched length
is \(u-v/\alpha\).
It exceeds \(\alpha v/(1-\alpha)\)
exactly when
\[
                       r<\frac{\alpha(1-\alpha)}
                                          {1-\alpha+\alpha^2}.
\]
This proves the source's crossing comparison.
Substituting the two specified data pairs
at \(\alpha=3/4\) gives each stated
rational endpoint.
\end{proof}
\end{annexitem}

\begin{annexitem}{annex:p15.neyman.accuracy.duality}{thm:p15-neyman-accuracy-duality}{False coverage, acceptance and integrated length}
\begin{proof}
The event \(t\in I(X)\) is precisely
\(X\in A_I(t)\), by definition.
Taking its probability under any true
parameter proves the identity, retaining
all nuisance values. If one acceptance
family minimizes this probability for
every false target and true law, and
its inversion is an admissible interval
system, that system satisfies exactly
the shortest-system definition.
For a lower half-line procedure,
\(t\) is included precisely when
\(\ell(X)\leq t\); the corresponding
upper version is equally direct.
This proves both one-sided statements.

The Lebesgue length of an interval is
the integral of its membership indicator.
For a measurable membership graph,
Tonelli's theorem for this nonnegative
function gives
\[
 E_{s,\eta}\operatorname{len}(I(X))
   =E_{s,\eta}\int_{\mathbb R}\mathbf1_{\{t\in I(X)\}}\,dt
   =\int_{\mathbb R}P_{s,\eta}\{t\in I(X)\}\,dt.
\]
Infinite lengths and expectations are
permitted. Thus pointwise dominance
over the candidate values integrates
to the asserted expected-length
dominance. No pointwise comparison
of the lengths at individual observations
was used in this identity.
\end{proof}
\end{annexitem}

\begin{annexitem}{annex:p15.neyman.uniform.shortest}{thm:p15-neyman-uniform-shortest}{The full cube-volume lower bound}
\begin{proof}
For candidate \(t>0\), let \(a_t(x)\in[0,1]\)
be any competing measurable acceptance rule
with \(E_t a_t=\alpha\).
On the cube \(D_t=(0,t)^n\), the
null density is \(t^{-n}\), so
\[
                             \int_{D_t}a_t(x)\,dx=\alpha t^n.
\]
If the true \(s\geq t\), its cube
contains \(D_t\); nonnegativity therefore
gives \(E_s a_t\geq\alpha(t/s)^n\).
The maximum-based acceptance shell
\(\{q_nt\leq L_n\leq t\}\) lies
inside \(D_t\), up to null boundaries,
and attains exactly this value.

If \(s<t\), the part of \(D_t\)
outside \(D_s\) has volume \(t^n-s^n\)
and \(a_t\leq1\). Therefore
\[
 \int_{D_s}a_t
     \geq \max\{0,\alpha t^n-(t^n-s^n)\}.
\]
After division by \(s^n\), the lower
bound is
\(\max\{0,1-(1-\alpha)(t/s)^n\}\).
The shell attains it: if \(s<q_nt\)
it has zero intersection with \(D_s\);
otherwise that intersection has
volume \(s^n-(q_nt)^n\).
This proves the minimum against every
true law simultaneously, including
all randomized competitors.
Its inversion is the already proved
nonempty exact-coverage maximum interval.
The three cases give the displayed
coverage function. Since \(t/s<1\)
in the first false-target case and
\(t/s>1\) in the second, the
probability is strictly below
\(\alpha\) whenever \(t\ne s\).
\end{proof}
\end{annexitem}

\begin{annexitem}{annex:p15.neyman.triangle.shortest}{thm:p15-neyman-triangle-shortest}{Every conditional overlap and the explicit false-coverage function}
\begin{proof}
Fix candidate \(t\), true scale \(s\),
and an allowed true nuisance \(\eta\).
The true density is constant in the
second coordinate on each fiber:
\[
                       h_{s,\eta}(x)=2/s^2-s\eta+3\eta x.
\]
It is nonnegative on \(0<x<s\).
A competitor must select length
\(\alpha(t-x)\) within \((0,t-x)\)
for almost every \(0<x<t\).

If \(s\geq t\), the true conditional
support \((0,s-x)\) contains the
entire candidate support for \(x<t\).
Every competitor contributes at least
that prescribed length; any acceptance
outside it adds nonnegative probability.
For \(t\leq x<s\), the candidate rule
has no required acceptance, so zero
is minimal. The upper-placement rule
uses exactly the prescribed lengths
and zero elsewhere, attaining all
these minima.

If \(s<t\), consider \(0<x<s\).
The largest length that can be placed
outside the smaller true support
\((0,s-x)\), but within \((0,t-x)\),
is \(t-s\). Hence the overlap is
at least
\[
                \max\{0,\alpha(t-x)-(t-s)\}
                 =\max\{0,s-(1-\alpha)t-\alpha x\}.
\]
The upper interval
\([(1-\alpha)(t-x),t-x]\)
attains this bound on each fiber.
Multiplying by the nonnegative
\(h_{s,\eta}(x)\) and integrating
proves optimality for every allowed
nuisance. Null sets in the fiber
constraint do not affect this
absolutely continuous integral.
This is the claimed conditional
comparison class.

When \(t\leq s\), the achieved
probability is
\[
        \alpha\int_0^t h_{s,\eta}(x)(t-x)\,dx
        =\alpha(t/s)^2
                    \left[1-\frac{s^3\eta}{2}(1-t/s)\right].
\]
When \(s<t<s/(1-\alpha)\), let
\(b=[s-(1-\alpha)t]/\alpha\in(0,s)\).
The overlap equals \(\alpha(b-x)\)
for \(0<x<b\) and zero elsewhere.
The same integration with \(b\)
in place of \(t\) gives the
stated formula with \(z=b/s\).
For \(t\geq s/(1-\alpha)\) every
overlap is zero.

To verify unbiasedness explicitly,
let \(H_\rho(z)=z^2[1-\rho(1-z)/2]\).
It has \(H_\rho(0)=0\),
\(H_\rho(1)=1\), and
\[
                          H_\rho'(z)=z[2-\rho+(3\rho/2)z].
\]
For \(0\leq\rho\leq2\) this
is strictly positive on \((0,1)\);
at \(\rho=2\) it is \(3z^2\).
For \(-1<\rho<0\), the bracket
is at least \(2+\rho/2>0\).
Thus \(H_\rho(z)<1\) for \(z<1\).
The false-target cases have
\(z<1\), while \(t=s\) has
\(z=1\). This proves every
claimed coverage comparison.
\end{proof}
\end{annexitem}

\begin{annexitem}{annex:p15.neyman.regular.obstruction}{thm:p15-neyman-regular-obstruction}{The two-sided derivative contradiction without a boundary shortcut}
\begin{proof}
Write \(P_0=P_t\) and \(p_0=p_t\)
only within this proof.
\(L^1\) differentiability gives
\[
                  P_{t+h}(B)=P_0(B)+h\int_Bg+o(|h|)
\]
for every measurable \(B\), with
the remainder bounded uniformly
by the \(L^1\) derivative error.
Integrating over the whole space
gives \(\int g=0\).
On \(Z=\{p_0=0\}\), for every
measurable \(B\subseteq Z\),
\(P_0(B)=0\) and
\(P_{t+h}(B)\geq0\).
The positive and negative \(h\)
limits force \(\int_Bg=0\).
Applying this to the positive
and negative parts shows
\(g=0\) a.e. on \(Z\).
Consequently the score
\(w=g/p_0\) on \(\{p_0>0\}\),
zero on \(Z\), is integrable
under \(P_0\), has mean zero
and has positive absolute mean.

Suppose a region \(A\) of
\(P_0\)-probability \(\alpha\)
minimized \(P_s(A)\) for all
nearby \(s\ne t\) on both sides.
For any other \(B\) of that
same probability,
\(P_{t+h}(B)-P_{t+h}(A)\geq0\).
The two-sided derivative and
equality at \(h=0\) force
\[
                              \int_Bg=\int_Ag.
\]
We show this is impossible.
The nonzero mean-zero score
has sets \(\{w\leq u\}\)
and \(\{w\geq v\}\) of
positive probability for some
\(u<0<v\).
Choose a positive \(\epsilon\)
smaller than their probabilities
and than \(\alpha,1-\alpha\).
Select subsets \(L,H\) in
these two sets with probability
exactly \(\epsilon\), and
select a disjoint set \(D\)
outside \(L\cup H\) with
probability \(\alpha-\epsilon\).
The selections are possible
after taking \(\epsilon\)
sufficiently small. In detail,
on any measurable set \(E\)
the function
\(r\mapsto P_0(E\cap\{x_1\leq r\})\)
is continuous from zero to
\(P_0(E)\), since a law with
Lebesgue density gives each
coordinate hyperplane probability
zero. The intermediate value
theorem gives subsets of
any required intermediate mass.

The regions \(D\cup L\)
and \(D\cup H\) each
have probability \(\alpha\),
but their score integrals
differ by at least
\(\epsilon(v-u)>0\).
This contradicts the equality
for all such \(B\).
The contradiction applies in
every two-sided neighborhood
of \(t\), proving the
precise obstruction. It uses
neither an unjustified
pointwise likelihood threshold
on an arbitrary boundary nor
nonzero derivatives only at
some unrelated parameter.
\end{proof}
\end{annexitem}

\begin{annexitem}{annex:p15.neyman.normal.one.sided}{thm:p15-neyman-normal-one-sided}{Both one-sided optima and actual nonexistence of a shortest system}
\begin{proof}
The standard normal law of
\cref{def:p15-normal-law,thm:p15-normal-normalization}
has positive continuous density,
total integral 1 and symmetry
about zero. Its cumulative
function is continuous and
strictly increasing from zero
to one, so \(z_\alpha\)
exists uniquely and
\(\Phi(-z_\alpha)=1-\alpha\).
The normal transformation theorem
gives \((X-\mu)/\sigma\)
that standard law.
Consequently
\[
 P_\mu(X-\sigma z_\alpha\leq\mu)=\alpha,\qquad
 P_\mu(\mu\leq X+\sigma z_\alpha)=\alpha.
\]

For a fixed candidate \(t\),
the lower interval has
acceptance region
\(A^-_t=(-\infty,t+\sigma z_\alpha]\).
For any true \(s>t\), the
ratio of the \(s\)-density
to the \(t\)-density is
strictly increasing in \(x\),
since it is
\[
 \exp\left(\frac{(s-t)(x-t)}{\sigma^2}
                         -\frac{(s-t)^2}{2\sigma^2}\right).
\]
The Neyman--Pearson comparison,
applied to rejection complements,
makes this lower acceptance
region uniquely minimizing
\(P_s(A)\) among regions
with \(P_t(A)=\alpha\),
up to null sets.
The equality set is a
single point and has
probability zero.
This proves the best lower
one-sided criterion for all
\(s>t\). For \(s<t\),
the likelihood ratio is
strictly decreasing, and
the same argument makes
\(A^+_t=[t-\sigma z_\alpha,\infty)\)
the unique minimizing
acceptance region. This
proves the upper version.

If an exact-\(\alpha\)
shortest interval system
existed, compare it with
these two valid one-sided
systems at a fixed candidate
\(t\). On the side \(s>t\)
its acceptance region would
have to tie the unique
minimum \(A^-_t\), and
on the side \(s<t\)
it would have to tie
\(A^+_t\).
These half-lines differ
on a set of positive
normal probability for
every \(0<\alpha<1\).
No region can equal
both a.e. This contradicts
the shortest-system requirement.
The example uses known
\(\sigma\); no assertion
about the distinct unknown-
variance Student optimum
is imported from it.
\end{proof}
\end{annexitem}

\begin{annexitem}{annex:p15.neyman.unbiased.derivatives}{thm:p15-neyman-unbiased-derivatives}{Global coverage peaks and the local curvature condition}
\begin{proof}
Fix a candidate \(t\) and
nuisance value, and let
\(F(s)=P_{s,\eta}(A_I(t))\).
Global unbiasedness says
that \(F(t)=\alpha\)
is a maximum. At an
interior point, ordinary
two-sided differentiation
therefore gives \(F'(t)=0\).
If the second derivative
exists, the symmetric
second difference
\([F(t+h)+F(t-h)-2F(t)]/h^2\)
is nonpositive and tends
to \(F''(t)\), so
\(F''(t)\leq0\).
The stated differentiation-
under-the-integral hypothesis
identifies these derivatives
with the two displayed
density integrals.
Comparing their second values
among the specified admissible
unbiased class is exactly
the source's local sharpness
criterion.

For the counterexample, put
\(v=s^2\in[0,9/25)\).
The Bernoulli success probability
is \(1/2-v+4v^2\).
Its minimum is \(7/16>0\)
at \(v=1/8\), and throughout
this range it is less
than 1; at the upper
endpoint its limiting value
is \(823/1250<1\).
Thus these are valid
probability laws.
The function has value
\(1/2\), first derivative
zero and second derivative
\(-2\) at \(s=0\).
But for \(1/2<|s|<3/5\),
\(-s^2+4s^4>0\), so
the acceptance probability
exceeds \(1/2\).
The derivative conditions
therefore do not replace
the global inequality.
Finally the explicit false-
coverage functions of the
uniform and triangular
systems were proved to
be at most \(\alpha\),
with equality at the
true target. They do
satisfy that global
unbiasedness property.
\end{proof}
\end{annexitem}
